\chapter{Course Map, Tools, and How to Learn}

\section{How a Message Moves Through a System}
Start with an order: after a customer pays, the shop needs to tell a later program, such as inventory, “the order is paid.” If the payment program calls inventory directly, the message may be lost while the receiver is down. A messaging system lets the sender hand off a message for storage, so the receiver can read it at its own pace. Apache Kafka is a distributed messaging and event-streaming platform; this book uses a runnable small teaching system to explain the underlying problems. No Kafka knowledge is assumed, and the course protocol is not presented as Kafka's official protocol.

This book calls one storable, readable piece of data a \textbf{message} or \textbf{record}. A \textbf{producer} sends records; a \textbf{broker} accepts requests and stores them; a \textbf{consumer} actively requests and processes them. A messaging system is software that transfers and stores records between these roles. A successful send request, a record being read, and successful business processing are three distinct events; later chapters store and inspect their states separately.

A \textbf{topic} is a named category for related messages, such as \pathcode{orders}. A topic can be split into ordered \textbf{partitions}; this course uses partitions so different records can be appended and read independently. A record's \textbf{key} can associate or route it, its \textbf{value} is its content, and its \textbf{timestamp} is a time value stored with the record. The timestamp does not replace the \textbf{offset}, which determines order within a partition. Each offset is the record's number within its own partition, starting at 0. Offsets from different partitions cannot be compared, and there is no single global order for the whole topic.
\begin{center}
\small
\begin{tabular}{lll}
\toprule
Topic & Partition & Record offsets in order \\
\midrule
\pathcode{orders} & 0 & 0, 1 \\
\pathcode{orders} & 1 & 0, 1 \\
\bottomrule
\end{tabular}
\end{center}
The two offset-0 entries are different records. This book guarantees order within each partition, not an order between the two partitions.

A record's lifecycle is: the producer chooses a partition; the broker appends the record to that partition's log; a consumer actively fetches it; the program runs a processing callback; and only after successful processing does it save its consumption progress. Reading does not delete the log record, so another consumer or a restarted consumer can read it again. Log retention determines when old records are deleted; consumption progress says where a consumer resumes. They are separate states. Chapter 4 explains the failure window between processing and saving progress.

\begin{center}
\small
\begin{tabular}{p{.22\textwidth}p{.68\textwidth}}
\toprule
Chapter and steps & Question answered \\
\midrule
Chapter 1, 1--4 & How do we encode, append, and read records, and identify a crash tail?\\
Chapter 2, 5--8 & How do indexes speed lookup, how are logs segmented, and how do head retention and tail truncation differ?\\
Chapter 3, 9--12 & How do we route to partitions, discover broker locations, and recover TCP bytes into complete requests?\\
Chapter 4, 13--16 & How do we batch sends, read records, save consumption progress, and understand replay after failure?\\
Chapter 5, 17--20 & How do consumers share partitions, and how are stale-member requests rejected?\\
Chapter 6, 21--24 & How do brokers replicate logs, and which records may ordinary consumers see?\\
Chapter 7, 25--28 & How do we elect a replica after failure, repair divergence, and admit a replica again?\\
Chapter 8, 29--31 & How do we observe cross-layer failures, identify repeated sends, and make local effects replay-safe?\\
\bottomrule
\end{tabular}
\end{center}

\section{What This Course Builds—and What It Does Not}
This project makes the path above runnable: partition logs on disk, TCP requests, producing and consuming, consumer offsets, a single-topic consumer group, and pull-based replication with controlled failover among local brokers. It is not an official Kafka client and does not use the Kafka wire protocol; do not point the official \method{KafkaProducer} or \method{KafkaConsumer} at this program.

The three Java source roots are part of the course design. \pathcode{provided/} contains data types, I/O, transport, serialization scaffolding, lifecycle management, and test fixtures, keeping the focus on algorithms rather than repeated setup. \pathcode{exercises/} is the only workspace students need to complete step by step. All step types and public signatures are visible in advance; methods not yet reached throw \method{ExerciseNotImplementedException} with the step number. \pathcode{reference/} contains an isolated complete implementation for comparison after trying the contract and for course verification. Its class names match the exercise sources, so student and reference source roots must never be compiled together. Tests reuse the same black-box contract in both modes; the student runtime classpath excludes reference classes.

This is not production-grade Kafka: it has no Kafka wire-protocol compatibility, transactions, compression, compaction, complete controller consensus, or controller HA. Chapter 8's limited idempotent-producer protocol and durable local effect ledger are neither Kafka production protocols nor end-to-end transactions. They use a trusted teaching control plane within one JVM; although multiple brokers have TCP servers, threads, and independent disk directories, their replicas still share that control plane. The course can demonstrate controlled failover, but cannot establish consensus under arbitrary network partitions or production-grade availability. Relevant steps explain the specific differences.

\section{Preparing Java 21+ and Course Commands}
The repository uses the Gradle Wrapper \coursecmd{./gradlew}; Gradle 9.8.1 runs on Java 17+, while Java compile, test, and demo tasks require a Java 21+ JDK and compile with \pathcode{--release 21}. Maven, Docker, and a running Kafka are not required. The scripts use Bash 3.2+, \pathcode{find}, \pathcode{awk}, \pathcode{curl}, and either \pathcode{sha256sum} or \pathcode{shasum}; \pathcode{tar} and \pathcode{install} prepare Tectonic, while macOS Pandoc extraction also needs \pathcode{unzip}. EPUB diagrams require Poppler's \pathcode{pdftoppm}; \coursecmd{./gradlew :setupBook} checks for it but does not install system packages. Install Poppler on Linux/WSL with \coursecmd{sudo apt-get install poppler-utils}, or on macOS with \coursecmd{brew install poppler}. First check \coursecmd{java -version}, \coursecmd{javac -version}, and \coursecmd{./gradlew --version}. You may use any suitable JDK through \pathcode{JAVA_HOME} or \pathcode{PATH}; SDKMAN is optional. The project does not silently change SDKMAN, system packages, or shell configuration.

\coursecmd{./gradlew :setup} resolves the pinned JUnit 1.11.4 dependency and verifies its checksum through Gradle; it does not store JUnit under \pathcode{.tools/}. \coursecmd{./gradlew :setupBook} prepares checksum-verified Tectonic 0.17.0 and Pandoc 3.12.1 under \pathcode{.tools/}, and checks Poppler without installing system software. \coursecmd{./gradlew :doctor} checks Gradle, Java, managed JUnit, and optional book tools. \coursecmd{./gradlew :compile} compiles the student skeleton and tests; successful compilation does not mean the exercises are implemented. The default course and list language is English; \coursecmd{COURSE_LANG=zh ./gradlew :list} displays Chinese step titles. \coursecmd{./gradlew :book} builds the English PDF and EPUB at \pathcode{build/book/en/simpleKafka.pdf} and \pathcode{build/book/en/simpleKafka.epub}; \coursecmd{COURSE_LANG=zh ./gradlew :book} builds the Chinese pair at \pathcode{build/book/zh/simpleKafka-zh.pdf} and \pathcode{build/book/zh/simpleKafka-zh.epub}. Each command also updates that language's checked-in PDF and EPUB under \pathcode{docs/book/dist/}. Readers can open those distributed files without Java or build tools. The book task starts the Wrapper JVM but does not compile Java or resolve JUnit. The first build needs a network connection for the TeX bundle; after warming both languages, use \coursecmd{BOOK_OFFLINE=1 ./gradlew :book --offline} and \coursecmd{COURSE_LANG=zh BOOK_OFFLINE=1 ./gradlew :book --offline}. Gradle \pathcode{--offline} controls Gradle dependencies; \pathcode{BOOK_OFFLINE=1} selects the cached TeX bundle for both Tectonic builds.

\section{Understanding Binary Data and File I/O}
A file stores bytes; it does not remember field boundaries for a program. The UTF-8 character “值” occupies 3 bytes, so its value length in the course example is 3, not its character count 1. An ordinary record has length=28+3=31 and occupies 4+31=35 bytes in total. An encoding must specify field order, integer width, and byte order. This course uses big-endian order: for example, the four bytes of integer 1 are \pathcode{00 00 00 01}, with the most significant byte first.

A \method{ByteBuffer} is a byte view with a capacity, position, and limit. Position is the next read/write location; limit is the end of the current readable/writable region. In write mode, \method{flip()} resets position to zero and sets limit to the position just written, so those bytes can then be read; flip does not clear the data. A buffer position is a byte cursor, while a record offset in later exercises is a logical number within a partition; they cannot be substituted. A decoder should advance only past bytes it has successfully consumed.

\method{FileChannel} can read or write at a specified file byte position and can request that written data be forced to storage. One \method{channel.read} may fill only part of a buffer, and one \method{write} may write only part of it; the provided \method{ChannelIO.readFully/writeFully} loops to handle this. Step 2 uses a lock to make concurrent appends to one log complete one at a time, so batches cannot interleave. The lock protects a short shared-state operation; do not hold it while waiting for network I/O. After each batch, \method{force(false)} is a simplified request to persist file data, not a mathematical guarantee against power loss, device caches, media damage, or filesystem semantics. CRC32C detects accidental changes to record bytes; it does not repair errors or prove that data came from a trusted source.

\section{How to Read Test Failures}
First read the failed Step test name and assertion, then return to the contract in this chapter. The first exception in a test report is usually more useful than the last part of the stack trace: an unimplemented method identifies its step and method; course errors should preserve \method{ErrorCode}; an unexpected IOException may be wrapped as \method{STORAGE_ERROR} with its cause preserved. Do not map errors to an empty list or a successful default value just to make tests green. A single-step diagnostic command runs only that step and does not mean its cumulative dependencies are satisfied; the cumulative command reevaluates every earlier step.
The course runner prints named tests grouped by step, per-step and overall counts, failure assertions with exact method navigation, and rerun commands. The report is always in English; legacy JUnit XML is saved under \pathcode{build/student/reports/} or \pathcode{build/reference/reports/}. These custom JavaExec tasks print plain text. Standard Gradle \pathcode{:exercises:test} and \pathcode{:reference:test} tasks support IDE debugging and exact \pathcode{--tests} selection; their XML/HTML reports use \pathcode{build/<mode>/gradle-test-results/} and \pathcode{build/<mode>/gradle-test-reports/}. The runner retains the \pathcode{NO_COLOR} and \pathcode{TERM=dumb} color gate.

\begin{center}
\small
\begin{tabular}{p{.48\textwidth}p{.43\textwidth}}
\toprule
Command & Meaning \\
\midrule
\coursecmd{./gradlew :list} & Show all 31 steps, their editable methods, and prerequisites.\\
\coursecmd{./gradlew :stepTest -Pstep=12} & Diagnose Step 12 only; useful for locating the current problem when it fails.\\
\coursecmd{./gradlew :test -Pstep=12} & Run cumulative tests from Step01Test through Step12Test.\\
\coursecmd{./gradlew :referenceTest -Pstep=31} & Run the complete 31-step contract against the isolated reference implementation.\\
\coursecmd{./gradlew :referenceDemo} & Run the completed reference TCP/failover scenario.\\
\coursecmd{./gradlew :referenceConsistencyDemo} & Run the real TCP and disk consistency scenarios for Chapter 8.\\
\bottomrule
\end{tabular}
\end{center}

Step numbers accept decimal values from 1 through 31, with optional leading zeroes (for example, \pathcode{-Pstep=1} or \pathcode{-Pstep=01}). Chapter boundaries are Steps 4, 8, 12, 16, 20, 24, 28, and 31: a chapter is complete only when cumulative tests through that boundary all pass. Passing \coursecmd{./gradlew :test -Pstep=8} means Steps 1–8 all pass; passing \coursecmd{./gradlew :stepTest -Pstep=8} means only the Step 8 diagnostic ran and does not authorize progression. Full-course tests must not put reference sources on the student classpath, and reference mode cannot substitute for exercise progress.

\section{Concepts Already Introduced and Where to Go Next}
Keep three ideas in mind: a topic names a category, its partitions are independently ordered logs, and each record's offset has meaning only within its own partition; sending, reading, and business processing are distinct events, and reading does not delete a record; log retention determines when records are deleted, while consumption progress determines where one consumer resumes. Every later step introduces its new terms and invariants before relying on them; no Kafka abbreviation is prior knowledge.

Use these questions to navigate the next steps: \stepref{1}{Record encoding and decoding} explains how a record becomes bytes with clear boundaries; \stepref{4}{Crash recovery} shows how to distinguish an unfinished file tail from complete damaged data; \stepref{14}{Pull-based consumer} demonstrates how a consumer's read position advances; \stepref{22}{ISR and high watermark} explains which records ordinary consumers may see when copies of a log have different progress; \stepref{25}{Clean election and leader epochs} shows how to reject requests from an old broker role after another broker takes over writing.
